\section{Методика эксперимента}
\label{sec:setup}
\subsection{Агент и сравнения}
Оценивается \texttt{zai/glm-5.3-flash} с Pi 0.85.1 и Harbor 0.22
на Terminal-Bench 2.0~\citep{terminalbench}, лимит --- 200 ходов.
Успех --- двоичный вердикт проверяющей программы задачи. Анализируемое
подмножество содержит 87 задач; две задачи, чьи окружения не удалось
инициализировать, исключены.

Приводятся четыре сравнения, каждое со своим контролем полной истории,
сопоставленным по задаче и повтору; контроли не объединяются.
\begin{itemize}\itemsep0pt
\item \emph{Состояние модели}: четыре условия, три повтора,
 $k\in\{0,2,6\}$.
\item \emph{Гибридное состояние}: восемь условий, два повтора. Варианты
 убирают записи неудач, добавляют недавнюю историю или сочетают
 способ обновления (через инструмент или в тексте ответа) со схемой
 (исходной или адаптированной).
\item \emph{Состав автоматических записей}: четыре набора и контроль,
 один повтор. Контроль и самые богатые наборы пересекаются с первым
 гибридным повтором, поэтому это не независимое воспроизведение.
\item \emph{Пилот окон}: 12 задач, шесть условий ($k\in\{0,2,6,12,\infty\}$
 и контроль), три повтора. Сохранились только итоговые сводки.
\end{itemize}
Схема адаптировалась на тех же задачах, а слабые наборы записей
отброшены после предварительных прогонов, поэтому сравнения поисковые
и внутривыборочные. Во всех реализациях $k$ --- заданная настройка
удержания; равное $k$ не гарантирует одинаковой сериализации.

\subsection{Единый набор показателей}
Каждое сравнение приводится в двух панелях с одинаковыми строками.
\emph{Панель качества}: доля решённых задач по всем включённым исходам,
медиана ходов $T$, медианная нормированная стоимость, медиана парного
отношения стоимости к полной истории и решения на доллар $\sum Y/\sum D$
(с учётом неудач). \emph{Панель ресурсов}: медианы пикового промпта,
накопленного некешированного входа $U$, чтения из кеша $K$, выхода $O$
и cache hit
\begin{equation}
 h_{iar}=\frac{K_{iar}}{U_{iar}+K_{iar}},
 \label{eq:cache-hit}
\end{equation}
который усредняется медианой по выполнениям, а не делением медиан
токенов. Счётчика записи в кеш нет. Нормированная стоимость
\begin{equation}
 D=10^{-6}(0.15U+0.03K+0.50O)
 \label{eq:referencecost}
\end{equation}
выражена в долларах США и не является счётом провайдера. Показатели,
которые сохранившиеся записи не позволяют восстановить (пики и cache hit
гибридов, доли решённых и входные счётчики пилота), обозначены --
и не подставляются.

\subsection{Статистический анализ и когорты}
Сопоставленные исходы проверяются двусторонним точным тестом Мак-Немара;
для состояния модели добавлены 95\% бутстрэп-интервалы по задачам
(100\,000 выборок, seed 20260922) с сохранением всех условий и повторов
задачи. Проверки поисковые, без поправки на множественность; незначимость
не означает эквивалентности. Когорты по горизонту задаются медианой ходов
контроля: короткая $T\leq8$ (31 задача), средняя $8<T\leq20$ (29),
длинная $T>20$ (27); состав фиксирован для всех условий и описывает
поведение контроля, а не назначенную сложность. Повторяемость в пилоте ---
число вызовов инструментов, делённое на число различных пар
инструмент/аргументы; 1.00 означает отсутствие идентичных повторов,
а не отсутствие смысловой избыточности. Доступные материалы перечислены
в дополнении~\ref{app:reproducibility}.
